\chapter{Conclusion}
\label{chap:conclusion}

This chapter concludes the thesis by summarizing the primary findings and contributions, discussing the limitations of our proposed approach, and outlining potential directions for future research.

\section{Summary of Contributions}
\label{sec:conclusion_summary}

In this thesis, we set out to address key structural weaknesses in multimodal vision-language models for document key information extraction, focusing specifically on LayoutLMv3. While pre-trained multimodal Transformers provide strong baseline representations by jointly processing text tokens, 2D coordinates, and visual image patches, their standard fine-tuning formulation treats spatial positions as flat numbers and classifies subwords largely in isolation. As observed in preliminary experiments, this causes frequent entity fragmentation and boundary confusion between adjacent text fields located on the same horizontal line.

To address these challenges without incurring the massive computational cost of pre-training a foundation model from scratch, we developed four targeted architectural enhancements:
\begin{enumerate}
    \item \textbf{Hierarchical Position Encoding (HPE):} We enriched token inputs with explicit line, column, and block structural indices derived from deterministic spatial grouping heuristics, allowing the backbone to distinguish reading units that share similar absolute coordinates.
    \item \textbf{Segment Context Module (SCM):} We introduced a segment-level Transformer layer stabilized by a ReZero residual gate and a first-token onset embedding. This mechanism pools subwords into cohesive word segments, aggregates inter-segment context, and broadcasts consistent entity information back to each individual token.
    \item \textbf{Intra-Line Boundary Loss:} We designed an auxiliary objective supervised by adjacent token pairs along the same text line. By directly penalizing false transitions between neighboring words, this loss encourages the model to learn sharp entity demarcations.
    \item \textbf{Cross-Modal Contrastive Loss:} We incorporated an InfoNCE alignment loss between text token embeddings and their corresponding image patch features, preventing visual drift during downstream fine-tuning.
\end{enumerate}

Comprehensive experimental evaluation on two established benchmarks---FUNSD (complex administrative forms) and CORD (commercial shopping receipts)---demonstrated the practical effectiveness of these additions across three random seeds:
\begin{itemize}
    \item On FUNSD, our enhanced model achieved an entity-level macro F1 score of $92.68\%$, improving upon the standard LayoutLMv3-base baseline ($90.31\%$) by $+2.37$ points.
    \item On CORD, the model attained $97.79\%$ F1, representing a $+1.28$ point gain over the base model ($96.51\%$) and performing competitively against much larger architectures.
    \item Detailed error analysis revealed a $36.6\%$ reduction in entity boundary errors and a $31.8\%$ drop in intra-line entity mergers, confirming that our design directly addresses the failure modes identified at the outset of this research.
\end{itemize}

\section{Limitations}
\label{sec:conclusion_limitations}

Despite the encouraging improvements, there are several practical limitations to our current implementation that should be acknowledged:
\begin{itemize}
    \item \textbf{Dependence on upstream OCR quality:} Our hierarchical grouping heuristics (line, column, and block clustering) rely on accurate word bounding boxes generated by the OCR engine. In scenarios where OCR text detection produces skewed, merged, or severely misaligned boxes, layout grouping errors can cascade into the hierarchical embeddings.
    \item \textbf{Heuristic assumptions for non-standard layouts:} The rule-based grouping logic assumes predominantly Manhattan-style layouts (horizontal text lines arranged in standard vertical columns). It may perform sub-optimally on documents with curved text, rotated stamps, diagonal receipts, or non-rectilinear artistic promotional flyers.
    \item \textbf{Training memory overhead:} While the auxiliary boundary classifier and contrastive projection heads are discarded during inference---resulting in zero additional runtime latency during deployment---they do increase GPU memory consumption by roughly $14\%$ during training. This necessitates careful batch size selection when training on hardware with limited VRAM.
\end{itemize}

\section{Future Directions}
\label{sec:conclusion_future_work}

Building upon the findings of this thesis, several promising directions remain for future exploration:
\begin{enumerate}
    \item \textbf{Learned layout grouping:} Replacing deterministic geometric thresholds with a lightweight graph neural network or clustering head could make hierarchical position extraction more robust to warped, rotated, and irregular document scans.
    \item \textbf{Knowledge distillation for edge deployment:} In many enterprise applications, invoices and receipts must be parsed locally on mobile devices or low-power servers. Distilling the predictions of our enhanced model into a compact student network (such as a 4-layer or 6-layer architecture) would be a valuable step toward practical real-time deployment.
    \item \textbf{Evaluation on Vietnamese documents:} An immediate and practical next step is adapting and evaluating this approach on Vietnamese document datasets, such as the MC-OCR dataset and domestic administrative forms. Because Vietnamese documents frequently contain multi-syllable compound words and distinctive diacritical marks, studying the interaction between syllable-level tokenization and segment-level context represents a compelling avenue of research.
\end{enumerate}
